\section{Proposed Method: LTD-Attack}\label{sec:method}
\LTD{} is organized around two complementary temporal scales. The \emph{micro} branch represents packet-level behavior within an individual page visit, whereas the \emph{macro} branch is intended to represent longitudinal characteristics across repeated observations. The hybrid combines their outputs to make a closed-world website prediction. This distinction follows the original research hypothesis, but this manuscript does not assume that the architecture described in the preliminary workshop summary is identical to the current frozen implementation.

\paragraph{Within-visit representation.} A micro encoder extracts information from individual packet sequences. Its concrete learned representation, sequence length, channel semantics, CNN/Transformer structure, and checkpoint provenance must be derived from the final implementation before submission. In particular, the earlier three-page abstract described a 3,000-packet CNN--Transformer, which is a historical description rather than verified documentation of the currently evaluated checkpoint.

\paragraph{Longitudinal representation.} The macro branch captures engineered descriptors or temporal summaries of repeated website observations. The earlier manuscript used 45 selected features and a heuristic fusion weight; neither number should be carried forward as an established property of the frozen model without a source-to-checkpoint audit. The crucial methodological questions are how a daily profile is constructed, what histories are allowed at prediction time, and whether class-conditioned aggregation could introduce information that would be unavailable to an actual observer.

\paragraph{Fusion and decision.} The currently evaluated model family includes Micro, Macro, and Hybrid variants. We deliberately use the abstract fusion notation in Eq.~\eqref{eq:hybrid} until the frozen training code, feature selection, XGBoost role, and score-combination recipe have been traced end-to-end. This is preferable to publishing a visually polished but inaccurate model diagram.

\paragraph{Diagram plan.} Figure~\ref{fig:research-pipeline} is an \emph{editable TikZ diagram of the experimental research workflow}, not yet a definitive diagram of the trained network internals. After code audit, a separate camera-ready Figure 1 will show exact encoder layers, longitudinal aggregation, fusion operations, and output head, with a common SVG source for future animation.
\begin{figure}[htbp]
\centering
\begin{tikzpicture}[>=Latex, node distance=5mm and 11mm,
  every node/.style={font=\small},
  stage/.style={draw=black!70,rounded corners,align=center,text width=3.2cm,minimum height=1.15cm,fill=black!2},
  eval/.style={draw=black!70,rounded corners,align=center,text width=3.0cm,minimum height=1.15cm,fill=black!8}]
\node[stage] (dev) {Historical DEV\\57,916 traces};
\node[stage,right=of dev] (b) {Chronological phase B\\Fit: Early+Middle\\Select: Late};
\node[stage,right=of b] (c) {Frozen phase C\\Refit: all Historical DEV};
\node[eval,below=12mm of c] (test) {Fixed evaluation\\09B / Future-B};
\node[eval,left=of test] (audit) {UID/class alignment\\Metrics and paired analysis};
\draw[->,thick] (dev)--(b);
\draw[->,thick] (b)--(c);
\draw[->,thick] (c)--(test);
\draw[->,thick] (test)--(audit);
\end{tikzpicture}
\caption{Editable diagram of the experimental pipeline, \emph{not} the final internal neural architecture. Phase B selection and phase C refit use only historical development data; 09B and Future-B are frozen retrospective evaluations.}
\label{fig:research-pipeline}
\end{figure}

